\documentclass[11pt]{article}
\usepackage[margin=1in]{geometry}
\usepackage{amsmath,amssymb,amsthm,mathtools}
\usepackage[hidelinks]{hyperref}

\newtheorem{theorem}{Theorem}[section]
\newtheorem{lemma}[theorem]{Lemma}
\newtheorem{proposition}[theorem]{Proposition}
\newtheorem{corollary}[theorem]{Corollary}
\newcommand{\E}{\mathbb E}
\newcommand{\Pp}{\mathbb P}
\newcommand{\Nn}{\mathbb N}
\newcommand{\cx}{\le_{\mathrm{cx}}}
\newcommand{\icx}{\le_{\mathrm{icx}}}

\title{Ballisticity for One-Dimensional Positive Cookie Environments Without Bounded Stacks}
\author{Working proof package for Kosygina--Zerner Problem 5.3}
\date{18 September 2026}

\begin{document}
\maketitle

\begin{abstract}
Kosygina and Zerner asked whether the bounded-stack assumption in the one-dimensional excited-random-walk ballisticity criterion can be replaced by positivity of every cookie.  We give a proof, conditional only on the standard monotonicity and backward-branching identities already used in the bounded theory.  The only new endpoint is $\delta=2$.  Bernoulliizing each positive cookie and moving the resulting deterministic-right cookies to the front reduces the endpoint to a critical Galton--Watson chain with random emigration.  Every integrable mean-two emigration law with positive mass at zero dominates, in the convex order relevant to the branching kernel, an explicit three-point bounded mean-two emigration law.  The bounded comparator is exactly a source-valid bounded cookie environment at $\delta=2$, so its stationary first moment is infinite.  Increasing-convex comparison transfers that divergence to the arbitrary random-emigration chain.  This yields zero speed at $\delta=2$; finite-coordinate strengthening gives all $\delta<2$, while finite-depth truncation gives positive speed for $\delta>2$.
\end{abstract}

\section{Source statement and notation}

We use the notation of Kosygina--Zerner, \emph{Excited random walks: results, methods, open problems}, arXiv:1204.1895v2, Problem 5.3.  In dimension one, at each site $x$ there is a cookie stack $(p_{x,i})_{i\ge1}$, where $p_{x,i}$ is the conditional probability of stepping to the right on the $i$th visit to $x$.  We assume:
\begin{itemize}
\item[(IID)] the site stacks are independent and identically distributed;
\item[(POS1)] $p_{x,i}\ge1/2$ almost surely for every $i$;
\item[(WEL)] the source weak-ellipticity condition: for each direction there is positive environment probability that all cookie transition probabilities in that direction are positive.
\end{itemize}
Set
\[
 d_i:=2p_{0,i}-1\in[0,1],
 \qquad
 \delta:=\E\sum_{i\ge1}d_i\in[0,\infty].
\]
The bounded-stack theorem states, under (IID), (BD), and (WEL), that the limiting speed $v$ satisfies $v=0$ for $\delta\in[-2,2]$ and $v>0$ for $\delta>2$.  Problem 5.3 asks whether (BD) may be replaced by (POS1).

Our target is therefore the following positive branch.

\begin{theorem}[Kosygina--Zerner Problem 5.3, positive branch]\label{thm:main}
Under \emph{(IID)+(WEL)+(POS1)},
\[
 \delta\le2 \Longrightarrow v=0,
 \qquad
 \delta>2 \Longrightarrow v>0.
\]
The second implication includes $\delta=+\infty$.
\end{theorem}

We use two standard facts from the one-dimensional cookie-walk theory.

\begin{enumerate}
\item \emph{Monotonicity.} Strengthening cookies coordinatewise cannot decrease the limiting speed.
\item \emph{Backward branching representation.} In the right-transient regime, let $V$ denote the usual backward branching-like chain.  When $V$ is positive recurrent with invariant law $\pi$, finite invariant mean gives
\[
 v=\frac{1}{1+2\E_\pi V};
\]
if $\E_\pi V=\infty$, the speed is zero.  These are the standard downcrossing/hitting-time identities used in the bounded proof.
\end{enumerate}

The only genuinely new part of the argument is to prove the infinite invariant first moment at the critical value $\delta=2$ without a bounded stack.

\section{The easy parameter regimes}

\subsection{The case $\delta>2$}

For $M\ge1$ set
\[
 \delta_M:=\E\sum_{i=1}^{M}d_i.
\]
By monotone convergence, $\delta_M\uparrow\delta$.  If $\delta>2$, including $\delta=+\infty$, there exists a finite $M$ with $\delta_M>2$.  Replace every cookie after depth $M$ by a fair cookie.  The resulting environment satisfies (IID)+(BD)+(WEL), has drift parameter $\delta_M>2$, and therefore has positive speed by the bounded theorem.  The original environment is coordinatewise stronger, so monotonicity gives $v>0$.

\subsection{Reduction of $\delta<2$ to $\delta=2$}

Suppose $\delta<2$.  Since $\sum_i\E d_i=\delta<\infty$, we have $\E d_i\to0$ and hence $\E(1-d_i)\to1$.  We may choose finitely many coefficients $a_i\in[0,1/2]$ so that
\[
 \sum_i a_i\E(1-d_i)=2-\delta.
\]
Define, for those finitely many coordinates,
\[
 p_i^+:=p_i+a_i(1-p_i),
\]
and leave all other coordinates unchanged.  The new drift satisfies
\[
 2p_i^+-1=d_i+a_i(1-d_i),
\]
so the strengthened environment has drift exactly two.  It remains (IID)+(POS1).  It also remains (WEL): if the original left probability $1-p_i$ is positive then the strengthened left probability $(1-a_i)(1-p_i)$ is still positive, and only finitely many coordinates were changed.  Thus, once the endpoint $\delta=2$ has been proved to have zero speed, monotonicity yields
\[
 0\le v\le v^+=0.
\]
It remains only to treat $\delta=2$.

\section{Bernoulliization and front-loading at the endpoint}

Assume from now on that $\delta=2$.

\subsection{Bernoulliization}

Conditionally on a site stack $(p_i)_{i\ge1}$, let
\[
 A_i\sim\mathrm{Bernoulli}(d_i),\qquad d_i=2p_i-1,
\]
independently in $i$, and let $C_i$ be independent fair Bernoulli variables.  Define the actual success bit
\[
 B_i=\begin{cases}
 1,&A_i=1,\\
 C_i,&A_i=0.
 \end{cases}
\]
Then
\[
 \Pp(B_i=1\mid(p_j)_j)=d_i+\frac{1-d_i}{2}=p_i,
\]
and the $B_i$ are conditionally independent.  Thus this enlargement leaves the annealed cookie coin sequence, and hence the annealed backward branching chain, unchanged.

Let
\[
 N:=\sum_{i\ge1}A_i.
\]
By Tonelli,
\[
 \E N=\E\sum_i d_i=2,
\]
so $N<\infty$ almost surely.

Moreover,
\begin{equation}\label{eq:p0}
 \Pp(N=0)>0.
\end{equation}
Indeed, the left-direction part of (WEL) gives positive probability of a stack on which $p_i<1$ for every $i$, equivalently $d_i<1$ for every $i$.  Since $\sum_i d_i<\infty$ almost surely, on that event
\[
 \prod_i(1-d_i)>0.
\]
This product is the conditional probability that every $A_i$ vanishes.  Averaging gives \eqref{eq:p0}.

\subsection{Front-loading}

Fix a realization of the marked/fair sequence.  Swap adjacent pairs of the form
\[
 [\text{fair bit},\,1]\quad\longmapsto\quad[1,\,\text{fair bit}].
\]
Each swap weakly increases the cumulative number of successes at every prefix.  After all swaps, the first $N$ cookies are deterministic successes and all later cookies are fair.  If $F_m$ is the number of failures before the $m$th success, then under this coupling
\begin{equation}\label{eq:front}
 F_m^{\mathrm{front}}\le F_m^{\mathrm{original}}
 \qquad\text{for every }m.
\end{equation}
Because $m\mapsto F_m$ is nondecreasing, the associated backward chains can be coupled in the same order.

For the front-loaded chain, conditionally on $N$, the failures after the deterministic successes are sums of iid geometric$(1/2)$ variables on $\{0,1,2,\dots\}$ with mean one.  Thus the endpoint is reduced to the following statement.

\begin{theorem}[Critical random-emigration theorem]\label{thm:RE}
Let $N$ be any nonnegative integer-valued random variable with
\[
 \E N=2,\qquad \Pp(N=0)>0.
\]
Let $N_n$ be iid copies of $N$, and let $\xi_{n,j}$ be iid variables with
\[
 \Pp(\xi=k)=2^{-(k+1)},\qquad k\ge0,
\]
independent of the $N_n$.  Define
\begin{equation}\label{eq:RE}
 Z_{n+1}=\sum_{j=1}^{(Z_n+1-N_n)_+}\xi_{n,j}.
\end{equation}
Then the positive recurrent chain \eqref{eq:RE} has invariant law $\pi_N$ satisfying
\[
 \E_{\pi_N}Z=\infty.
\]
\end{theorem}

Positive recurrence in Theorem~\ref{thm:RE} is elementary from the bounded theory: since $\E(N\wedge M)\uparrow2$, choose $M$ with $\E(N\wedge M)>1$.  The chain with bounded emigration $N\wedge M$ is the backward chain of a bounded positive cookie environment with drift greater than one and is positive recurrent.  The chain with emigration $N$ is pathwise no larger, so its return time to zero is no larger.

\section{Convex-order proof of the critical random-emigration theorem}

\subsection{A universal three-point comparator}

Set
\[
 p_0:=\Pp(N=0),
 \qquad
 a:=\min\left\{\frac{p_0}{2},\frac16\right\}>0.
\]
Define $N_*$ by
\begin{equation}\label{eq:Nstar}
 \Pp(N_*=0)=a,
 \quad
 \Pp(N_*=2)=1-3a,
 \quad
 \Pp(N_*=3)=2a.
\end{equation}
Then $\E N_*=2$.

\begin{lemma}\label{lem:cx}
We have
\[
 N_*\cx N.
\]
\end{lemma}

\begin{proof}
For integrable integer-valued variables of equal mean, convex order is equivalent to the stop-loss inequalities
\[
 \E(X-k)_+\le\E(Y-k)_+\qquad(k=0,1,2,\dots).
\]
This is the discrete hinge representation of a convex function: its second differences are nonnegative, so after subtracting an affine function it is a nonnegative linear combination of $(x-k)_+$.

At $k=0$ both stop-loss transforms equal the common mean $2$.  At $k=1$,
\[
 \E(N-1)_+=\E N-\Pp(N\ge1)=1+p_0,
\]
whereas
\[
 \E(N_*-1)_+=(1-3a)+2(2a)=1+a\le1+p_0.
\]
At $k=2$, equality of the mean to two gives
\[
 \E(N-2)_+=\E(2-N)_+=2p_0+\Pp(N=1)\ge2a,
\]
while
\[
 \E(N_*-2)_+=2a.
\]
For $k\ge3$ the stop-loss transform of $N_*$ is zero.  This proves the claim.  (For real thresholds, both stop-loss transforms are affine between consecutive integers.)
\end{proof}

\subsection{The random-sum kernel preserves increasing convexity}

Let
\[
 S_r:=\xi_1+\cdots+\xi_r,\qquad S_0=0.
\]
For an increasing convex function $\varphi:\mathbb Z_+\to[0,\infty)$ put
\[
 h(r):=\E\varphi(S_r).
\]

\begin{lemma}\label{lem:hconvex}
The function $h$ is increasing and convex on $\mathbb Z_+$.
\end{lemma}

\begin{proof}
It is increasing because $S_{r+1}=S_r+\xi_{r+1}\ge S_r$ under the natural coupling.  For convexity define
\[
 g(s):=\E[\varphi(s+\xi)-\varphi(s)].
\]
For each fixed $k$, convexity of $\varphi$ makes $s\mapsto\varphi(s+k)-\varphi(s)$ nondecreasing; hence $g$ is nondecreasing.  Therefore
\[
 h(r+1)-h(r)=\E g(S_r)
\]
is nondecreasing in $r$, because $S_{r+1}\ge S_r$ stochastically.
\end{proof}

For an emigration law $L$, write
\[
 (K_L\varphi)(x):=\E\varphi\bigl(S_{(x+1-L)_+}\bigr).
\]

\begin{lemma}\label{lem:kernel}
For every increasing convex $\varphi$,
\[
 K_{N_*}\varphi(x)\le K_N\varphi(x)\qquad(x\ge0),
\]
and, for either emigration law, $x\mapsto K_L\varphi(x)$ is increasing and convex.
\end{lemma}

\begin{proof}
Extend the discrete convex function $h$ from Lemma~\ref{lem:hconvex} linearly between consecutive integers; the extension is convex and nondecreasing.  For fixed $x$, the function
\[
 q_x(n):=(x+1-n)_+
\]
is convex in $n$.  Hence $n\mapsto h(q_x(n))$ is convex.  Lemma~\ref{lem:cx} yields the first inequality.

For the second assertion, fix $l$.  The map $x\mapsto(x+1-l)_+$ is increasing and convex; composing it with the convex nondecreasing $h$ remains increasing and convex.  Averaging over $l$ preserves these properties.
\end{proof}

\begin{proposition}\label{prop:icx}
Let $D$ be the chain \eqref{eq:RE} with emigration $N_*$, and let $Z$ be the chain with emigration $N$.  If $D_0=Z_0=0$, then
\[
 D_n\icx Z_n\qquad\text{for every }n\ge0.
\]
In particular,
\[
 \E_0D_n\le\E_0Z_n.
\]
\end{proposition}

\begin{proof}
The assertion is trivial at time zero.  If it holds at time $n$, then for any increasing convex $\varphi$, Lemma~\ref{lem:kernel} gives
\[
\begin{aligned}
 \E\varphi(D_{n+1})
 &=\E(K_{N_*}\varphi)(D_n)\\
 &\le\E(K_N\varphi)(D_n)\\
 &\le\E(K_N\varphi)(Z_n)\\
 &=\E\varphi(Z_{n+1}).
\end{aligned}
\]
The middle inequality uses that $K_N\varphi$ is increasing convex.  Induction completes the proof.
\end{proof}

\subsection{The comparator has divergent stationary mean}

The law $N_*$ corresponds exactly to a bounded iid cookie environment: at a site, the first $N_*$ cookies are deterministic-right and every later cookie is fair.  This environment satisfies (BD) with bound three and has drift parameter
\[
 \delta_*=\E N_*=2.
\]
It also satisfies the source (WEL): every right probability is positive; for the left direction, the event $\{N_*=0\}$ has probability $a>0$ and makes all cookies fair.

Consequently the bounded theorem gives speed zero.  Its backward chain is exactly $D$.  Since $\delta_*=2>1$, $D$ is positive recurrent; if its invariant mean were finite, the standard speed identity would give strictly positive speed.  Hence, with invariant law $\pi_*$,
\begin{equation}\label{eq:infmeanD}
 \E_{\pi_*}D=\infty.
\end{equation}

The chain is irreducible and aperiodic: from zero, $N_*=0$ followed by one geometric variable reaches every state with positive probability, and $0$ has a self-loop.  Therefore
\[
 \mathcal L_0(D_n)\longrightarrow\pi_*.
\]
For every finite $R$,
\[
 \E_0(D_n\wedge R)\longrightarrow\E_{\pi_*}(D\wedge R).
\]
Using \eqref{eq:infmeanD} and then letting $R$ grow shows
\begin{equation}\label{eq:EDdiverge}
 \E_0D_n\longrightarrow\infty.
\end{equation}

\subsection{Completion of Theorem~\ref{thm:RE}}

The kernel of $Z$ is stochastically monotone in the state: for $x\le y$, use the same $N$ and the same geometric sequence to obtain
\[
 \sum_{j=1}^{(x+1-N)_+}\xi_j
 \le
 \sum_{j=1}^{(y+1-N)_+}\xi_j.
\]
Let $\pi_N$ be the invariant law.  Suppose, toward a contradiction, that
\[
 \E_{\pi_N}Z<\infty.
\]
Couple a chain $Z^{(0)}$ started at zero with a stationary chain $Z^{(\pi)}$ started with law $\pi_N$, using identical transition randomness.  Since $0\le Z^{(\pi)}_0$,
\[
 Z_n^{(0)}\le Z_n^{(\pi)}\qquad\text{for all }n.
\]
Thus
\[
 \E_0 Z_n\le \E_{\pi_N}Z
\]
uniformly in $n$.  On the other hand Proposition~\ref{prop:icx} and \eqref{eq:EDdiverge} give
\[
 \E_0 Z_n\ge\E_0D_n\longrightarrow\infty,
\]
a contradiction.  Therefore $\E_{\pi_N}Z=\infty$, proving Theorem~\ref{thm:RE}.

\section{Return to the original cookie environment}

By front-loading, the random-emigration chain of Theorem~\ref{thm:RE} is pathwise no larger than the Bernoulliized backward chain of the original critical environment.  The front chain has infinite invariant first moment.  Both chains are positive recurrent; using the same finite-time/stationary comparison as above forces the original critical backward chain to have infinite invariant first moment as well.  Hence its speed is zero.  This proves the endpoint $\delta=2$.

Section 2 then gives every $\delta<2$ case by finite-coordinate strengthening and every $\delta>2$ case by finite-depth truncation.  Theorem~\ref{thm:main} follows.

\section{Prior-art and verification status}

This manuscript is a working proof package, not a novelty assertion.  A dated search on 18 September 2026 found the following especially relevant works.

\begin{itemize}
\item O. Chakhtoun, \emph{One-Dimensional Excited Random Walk with Unboundedly Many Excitations Per Site}, Ph.D. dissertation, CUNY Graduate Center, 2019.  Its public abstract states that a tail-decay estimate on the number of excitations per site is assumed and that ballisticity/zero-speed criteria are obtained under that wider class.  The complete thesis must be compared theorem-by-theorem with the no-tail-assumption statement here.
\item T.-M. Nguyen, \emph{Speed of Excited Random Walks with Long Backward Steps}, J. Stat. Phys. (2022), which contains a critical multi-type Galton--Watson theorem with deterministic emigration for a different ERW model.
\item P. Carper, \emph{A Branching Process with Geometric Emigration and the ``Have Your Cookie and Eat It'' Random Walk}, arXiv:2609.06311 (submitted 5 September 2026), which treats a special geometric-emigration law.
\item Classical critical branching-with-emigration/migration papers of Vatutin, Kaverin, Nagaev--Khan, and Yanev--Mitov.  Their exact hypotheses should be checked before any novelty claim for Theorem~\ref{thm:RE}.
\end{itemize}

No bounded search located a later source explicitly claiming the full no-tail-assumption resolution of Kosygina--Zerner Problem 5.3.  That is negative evidence only.

Under the repository owner protocol this package remains operationally \textbf{PARTIAL} until it receives an actual independent P26 full-proof audit, one distinct OTHER-track audit, and complete prior-art reconciliation.  No finite computational core is used in the proof.

\begin{thebibliography}{9}
\bibitem{KZ}
E. Kosygina and M. P. W. Zerner,
\emph{Excited random walks: results, methods, open problems},
Bull. Inst. Math. Acad. Sin. (N.S.) 8 (2013), 105--157;
arXiv:1204.1895v2.

\bibitem{Chakhtoun}
O. Chakhtoun,
\emph{One-Dimensional Excited Random Walk with Unboundedly Many Excitations Per Site},
Ph.D. dissertation, CUNY Graduate Center, 2019.

\bibitem{Nguyen}
T.-M. Nguyen,
\emph{Speed of Excited Random Walks with Long Backward Steps},
J. Stat. Phys. (2022).

\bibitem{Carper}
P. Carper,
\emph{A Branching Process with Geometric Emigration and the ``Have Your Cookie and Eat It'' Random Walk},
arXiv:2609.06311, 2026.
\end{thebibliography}

\end{document}
